% !TeX root = ../../../main.tex
\section{Discussion}

\subsection{Clinical Observations}

This prospective case-control study of 64 renal transplant patients investigated the utility of semiquantitative metrics derived from VisR imaging for assessing renal allograft pathology. Tables~\ref{tbl:demographics}, \ref{tbl:biopsy}, and \ref{tbl:banffScores} summarize patient demographics, biopsy justifications and findings, and Banff scores for biopsies performed during this study, respectively. Overall, the patient cohort recruited for this study reflects a relatively healthy population, with few patients experiencing delayed graft function or graft rejection. However, since UNC does not perform protocol biopsies, this also means that VisR imaging data could only be matched to biopsy findings for zero-hour biopsies and indication biopsies. As such, many VisR measurements could not be associated with biopsy findings due to timing mismatches between imaging and biopsy dates. Fig.~\ref{fig:ptTime} also shows that the availability of VisR imaging data over time greatly varies by patient, with some patients only undergoing imaging on a single date while others had multiple imaging sessions over several months. 

Most patient demographic information including age, sex, race, BMI, donor type, and reasons for study withdrawal were not significantly correlated with one another, with a few exceptions noted in the Results section. Importantly, many of the correlations (e.g. Asian or Native American racial identities, study departure due to pregnancy) can be attributed to small sample sizes in those demographic groups. However, the relationships concerning the living status of graft donors and White and Black racial identities may reflect underlying socioeconomic factors that influence access to living donors in these populations. Still, many of the correlation coefficients were relatively small ($|\rho| < 0.4$), indicating that these relationships were not particularly strong.

Fig.~\ref{fig:clinicalVisrExample} demonstrates that PD, RE, and RV images, on their own, qualitatively reflect differences in the viscoelasticity of renal allograft microstructure. The RE and RV images, in particular, denote the anisotropic nature of the renal parenchyma: the 0 degree orientation shows higher RE and RV distributions in the medulla compared to the cortex (and vice versa for PD), whereas the 90 degree orientation shows relatively homogeneous distributions except near the corticomedullary junction. It should be noted that, since the expected elastic contrast between the renal cortex or medulla in healthy kidneys is not fully established as discussed in Section~\ref{sec:about_offaxis}, these qualitative observations cannot be directly compared to existing literature. Nevertheless, these findings qualitatively agree with previous \emph{ex vivo} and \emph{in vivo} studies that have demonstrated the anisotropic viscoelastic nature of renal parenchyma. Furthermore, the ability to clearly distinguish between the renal cortex and medulla using PD and VisR parametric images are encouraging, as they demonstrate VisR's inherent capability to image renal mechanical information without the need for shear wave tracking and spatial averaging.

Recall that the positions of ROIs were determined using B-mode images with overlays of ARFI PD and log(VoA) images to ensure that ROIs were placed in the appropriate anatomical regions of the renal parenchyma. While it is known that B-mode images do not consistently distinguish anatomical regions or pathologies within the renal parenchyma~\cite{rodgers_Ultrasonographic_2014,cloutier_Quantitative_2021}, the addition of elastographic information helped to identify regions of differing stiffness that corresponded to the outer cortex, inner cortex, outer medulla, and inner medulla. Here, PD served as a reference viscoelasticity metric taken using an on-axis ARF-based method while log(VoA) provided complementary information regarding fluid-filled tissues where ARF-induced decorrelation can be expected~\cite{torres_Delineation_2019,anand_Comparing_2023,qazi_Comparison_2025}. While this has aided the placement of ROIs in this work as an ad-hoc solution, future work should characterize the utility of this approach in a systematic manner. Alternatively, log(VoA) itself may serve as a useful biomarker for renal allograft pathology.

\subsection{VisR DoA Findings by Biopsy Stratification}

% DoA box plots:
%
% For i:
% - Small sample sizes for i=2 and i=3 limit conclusions
%   since we had to group those with i>0 together
% - i>0 vs i=0 difference most noticeable in outer medulla
%   using PD, RE, and RV
% - i=1 vs i=0 appear different in outer cortex for VisR RE/RV,
%   but not to a significant amount; trend unclear for PD
% - Direction of DoA change in outer cortex vs outer medulla
%   are opposite when i>0; DoA more positive in outer cortex,
%   more negative in outer medulla
%   - If true, implies elasticity/viscosity increase along nephrons
%     and/or elasticity/viscosity decrease across nephrons in outer cortex
%     given the preesence of interstitial inflammation
%   - ...as well as elasticity/viscosity decrease along nephrons
%     and/or elasticity/viscosity increase across nephrons in outer medulla
In Fig.~\ref{fig:clinicalVisrDoa}, DoA distributions derived from VisR measurements exhibited significant differences across groups stratified by \emph{i} scores, IFTA grading, \emph{i-IFTA} scores, and findings of graft rejection in multiple anatomical regions of the renal parenchyma. Starting with the top row of the figure, DoA distributions calculated using PD, RE, and RV were significantly different between those associated with interstitial inflammation (\emph{i}) scores of zero from those with scores of one or greater in the outer medulla. Although the number of patients with \emph{i} scores of 2 and 3 were limited such that they had to be grouped together with \emph{i} score of 1 for analysis, these findings hint at the potential utility for VisR-derived DoA metrics to detect interstitial inflammation in renal allografts. It should be noted that the direction of change in DoA for PD with the presence of inflammation is opposite to that of RE- and RV-based DoA. This is expected; drastic viscous changes notwithstanding, lower PD values correspond to increased stiffness in underlying tissue. Together, this may indicate that interstitial inflammation is associated with relative increases in elasticity and viscosity along the dominant direction of nephrons (i.e. the longitudinal shear elastic modulus and viscosity) in the outer cortex, whereas the reverse is true in the outer medulla.

% For IFTA:
% - IFTA3 most clearly distinguished from IFTA<3 in outer and inner cortex
%   using PD, RE, and RV
%   - If true, implies elasticity/viscosity decrease along nephrons
%     and/or elasticity/viscosity increase across nephrons
%     when severe IFTA is present
% - RE, RV of IFTA1 and IFTA2 appear to fall between IFTA0 and IFTA3
%   in inner cortex; not distinguishable elsewhere or with PD
Continuing down to the second row of Fig.~\ref{fig:clinicalVisrDoa}, DoA distributions derived from PD, RE, and RV were significantly different between those without IFTA and those with grades of IFTA3 in the outer medulla, as well as between IFTA2 and IFTA3 in the outer cortex. Both changes show marked decreases in RE- and RV-based DoA compared to biopsies with less severe or no IFTA, whereas the opposite trend is again observed for PD-based DoA. This step in DoA was also observed for RE and RV-based DoA, but not PD, between IFTA0 and IFTA1 in the outer medulla. Together, these findings suggest that renal allografts with severe IFTA exhibit relative decreases in elasticity and viscosity along the dominant direction of nephrons. This trend appears counterintuitive; one may expect an increase in fibrosis to be associated with increased stiffness along the dominant direction of nephrons, instead. However, since the number of biopsied patients with high IFTA scores was limited in this study, future work with larger cohorts is needed to validate this finding.

% For i-IFTA:
% - i-IFTA3 not available due to no biopsies with that score
% - i-IFTA1/2 not significantly different in most regions, PD/RE/RV
% - i-IFTA2 vs 0 different in outer cortex, outer medulla w/ PD only
%   but not PD; i-IFTA2 not significantly different from either 0 or 1
%   regardless of metric
\emph{i-IFTA} scores, which reflect the degree of interstitial inflammation present in areas affected by IFTA, are used to classify DoA distributions in the third row of Fig.~\ref{fig:clinicalVisrDoa}. Although PD-based DoA distributions were significantly different between \emph{i-IFTA} scores of 0 and 2 in the outer cortex and outer medulla, RE and RV-based DoA distributions were not found to be significantly different across \emph{i-IFTA} scores anywhere in the renal parenchyma. In addition, no significant differences were observed between \emph{i-IFTA} scores of 1 and 2 in any anatomical region using any of the three metrics. These findings suggest that VisR-derived DoA metrics may be sensitive to mild \emph{i-IFTA} presence in renal allografts to a lesser extent than interstitial inflammation or IFTA alone. However, the limited sample sizes in this study, particularly for higher \emph{i-IFTA} scores, warrant further investigation with larger cohorts to validate these findings.

% For rejection:
% - no significant differences
% - RE, RV not significantly different but median appears to
%   follow opposite trend as PD (i.e. lower PD DoA but higher RE/RV DoA
%   in rejection vs non-rejection, similar to i and IFTA findings)
Regarding graft rejection, the bottom row of Fig.~\ref{fig:clinicalVisrDoa} shows that no significant differences were observed between DoA distributions associated with biopsies with findings of transplant rejection versus all other biopsy results. The aforementioned issue of small sample sizes is most noticeable here, as only four patients experienced graft rejection during the study period. It should be noted, however, that the Banff system does not exclusively attribute specific lesion types to rejection; rather, it considers the overall constellation of findings in a biopsy to determine whether rejection is present. As such, the "Not Rejected" category necessarily does not make a distinction regarding the presence of other rejection-relevant pathologies such as i-IFTA.

% Therefore:
% - Inflammation have most influence on DoA in outer cortex and medullae
%   using RE/RV; PD may also be useful in outer medulla only
% - IFTA findings have most influence on DoA in inner cortex
%   using PD, RE, and RV for severe IFTA
% - IFTA findings different in outer medulla for detecting IFTA presence
%   using PD, RE, and RV (but not useful for grading IFTA)
% - Rejection may be best detected in outer cortex only using PD
% - RE and RV being particularly similar; PD showed more differences
% - Sample size remains a major limitation to these analyses
Putting these findings together, it appears that renal allografts with interstitial inflammation, fibrosis, and the combination thereof exhibit changes in VisR-derived DoA metrics that differ between various anatomical regions of the renal parenchyma. The clearest changes in DoA were observed for IFTA3 versus lower IFTA grades using PD, RE, or RV in all layers up to the outer medulla, suggesting that severe IFTA may have a more pronounced effect on VisR-derived DoA metrics. 

\subsection{VisR RR Findings by Biopsy Stratification}

% RR boxplots:
% 
% For i:
% - Small sample sizes for i=2 and i=3 hurt us here, too
% - i>0 vs i=0 difference most noticeable in RR between
%   outer vs inner medulla
%   - ...but even then, difference is minimal for all metrics
Fig.~\ref{fig:clinicalVisrRr} shows RR distributions derived from VisR measurements that exhibited significant differences across groups stratified by \emph{i} scores, IFTA grading, \emph{i-IFTA} scores, and findings of graft rejection in multiple anatomical regions of the renal parenchyma. Starting with the top row of the figure, RR distributions calculated using PD were significantly different between those associated with interstitial inflammation (\emph{i}) scores of 0 from those with scores of 1 or greater in the RR between the outer and inner medulla. However, this difference appears to be minimal; although the nonparametric Kruskal-Wallis test yielded a significant $p$-value, the median and interquartile ranges of RR distributions for both groups appear to overlap considerably. No significant differences were observed using RE or RV-based RR distributions in any anatomical region for \emph{i} scores, though the results are confounded by the amalgamation of all \emph{i} scores greater than 0 into a single group.

% For IFTA:
% - RE/RV outer cortex vs outer medulla: IFTA2 vs 3 different (significant drop)
% - PD/RE/RV outer vs inner medulla: IFTA0 vs IFTA2 (drop more modest)
% - therefore, potential to classify IFTA0, IFTA3 using RR?
%
% For i-IFTA:
% - Only meaningful differences seen between:
%   - i-IFTA0 vs 1, 1 vs 2 in outer cortex vs outer medulla for PD only
%   - i-IFTA0 vs 2 for PD, RE, RR in outer medulla vs inner medulla 
%     - RR seems to falls over i-IFTA severity for RE/RV, rise for PD
%       (but spread of data too wide to make definitive claims)
However, the second row of Fig.~\ref{fig:clinicalVisrRr} shows that PD, RE, and RV-based RR distributions are significantly and noticeably different across IFTA severity. Namely, RE and RV-based RR distributions between the outer cortex and outer medulla were significantly different between IFTA2 and IFTA3, with a marked decrease in RR for IFTA3-associated biopsies. Likewise, PD, RE, and RV-based RR distributions between the outer and inner medulla were significantly different between IFTA0 and IFTA2, and the overall distribution of RE- and RV-based RR distributions appeared to decrease with increasing IFTA severity. However, this difference is very modest, partly owing to the limited sample sizes for higher IFTA grades in this study. Together, these findings suggest that VisR-derived RR metrics could be useful for assessing IFTA severity in renal allografts, particularly for distinguishing between no IFTA and severe IFTA. Likewise, the third row of Fig.~\ref{fig:clinicalVisrRr} shows that RR distributions derived from PD, RE, and RV were significantly different between \emph{i-IFTA} scores of 0 and 2 in the RR between the outer and inner medulla. Here, RE and RV-based RR values appeared to decrease with increasing \emph{i-IFTA} severity, whereas PD-based RR values appeared to increase with increasing \emph{i-IFTA} severity.

% For rejection:
% - Significant differences exist for all shown RRs except inner cortex
%   vs outer medulla regardless of metric used, outer cortex vs outer
%   medulla for PD only
% - RE, RV-based RR appear more sensitive to rejection-related changes
%   than PD-based RR given the lower p-values observed
Finally, the bottom row of Fig.~\ref{fig:clinicalVisrRr} shows that RR distributions derived from RE and RV were significantly different between biopsies indicating rejection and those that did not in all anatomical region pairs except for the inner cortex versus outer medulla. For all RE- and RV-based RR distributions where significant differences were observed, RR values appeared to decrease in rejected grafts compared to non-rejected grafts. This trend is similar to that observed for IFTA severity, particularly in the RR between the outer cortex and outer medulla. However, the connection between the two findings is unclear due to the limited number of biopsied and rejected patients in this study. 

% Therefore...
Putting these findings together, it appears that renal allografts with interstitial inflammation, IFTA, and graft rejection also exhibit changes in VisR-derived RR metrics. However, the clearest changes were observed for IFTA severity and graft rejection using RE and RV-based RR metrics, implying that these metrics are potentially sensitive to pathological changes in renal allografts than PD-based RR metrics. 

\subsection{Relating Pathologies to VisR DoA and RR Metrics}

Table~\ref{tbl:clinicalRegressor} demonstrated that linear and multinomial regression models could be developed to directly relate VisR-derived DoA and RR metrics to UPCR, IFTA grading, and \emph{i-IFTA} scores. However, while the models' accuracy metrics were significantly different from naive predictors of the mean of clinical variables (i.e. those that predict the root-mean-square of differences from the mean or return a lower cross-entropy loss), their actual performances were modest at best. This is most evident from the relatively low Spearman's rank correlation coefficients (i.e. average $|\rho_{s}|$ never exceeding 0.4 across all cross-validation folds) observed for all clinical variables. Cross-entropy loss and RMSE showed similarly modest improvements over naive predictors where significant differences arose. However, this does not necessarily preclude the clinical utility of developing VisR-based regression models for predicting clinically relevant metrics; it simply suggests that additional multicolinear factors may need to be considered to improve model performance. 

Likewise, Table~\ref{tbl:clinicalClassifier} demonstrated that VisR-derived DoA and RR metrics, overall, exhibited weak trends that support the development of simple predictive models based on the categories listed in Table~\ref{tbl:clinicalVariables}. Only four groups of models exhibited AUC distributions that were significantly and consistently greater than 0.5, all of which arose from IFTA classifiers involving the medulla. These include the multinomial logistic regression models using (1) RE- and (2) RV-based DoA in the outer medulla evaluated against IFTA1, as well as (3) PD-based DoA in the inner medulla and (4) PD-based RR between the outer and inner medulla evaluated against IFTA2. This is consistent with the findings from Fig.~\ref{fig:clinicalVisrDoa} and Fig.~\ref{fig:clinicalVisrRr}, which showed that VisR-derived DoA and RR metrics were most sensitive to IFTA severity in the medulla. It is notable that this finding does not extend to IFTA3 classification, though, likely due to the limited number of biopsies with that diagnosis grade in this study.

It should also be pointed out that Table~\ref{tbl:clinicalClassifier} shows many cases where the average specificity and sensitivity of models were significantly greater than 0.5, even when AUC distributions were not. This suggests that the models were able to correctly classify true negatives or true positives, respectively, at rates greater than random guessing at its optimal performance point even if their overall performances were inconsistent. For example, since multiple related decision boundaries adjacent were simultaneously created for the multinomial classifiers for IFTA and \emph{i-IFTA} (e.g. IFTA0 versus IFTA1 or above, IFTA1 versus IFTA2 or above, etc.), the AUC in each class may have been diluted in favor of achieving an improved balance overall. Still, classifiers that are biased towards either specificity or sensitivity may still have clinical utility depending on the application context. For instance, if VisR is positioned as a screening or tool for potential indication biopsies, ROC operating points that favor sensitivity over specificity may be appropriate. Likewise, VisR-based diagnostic pipelines that prioritize specificity may be useful for ruling out the need for biopsies in low-risk situations. The formal identification of decision thresholds and operating points, however, is deferred to future work.

%\subsection{Early Detection of Renal Allograft Pathologies}
%
% not worth it due to data quality (likely due to sample size issues)

\subsection{Limitations and Future Work}
\label{sec:transplantDiscussionLimitations}

As mentioned extensively throughout this section, the fundamental limitation of this study is the relatively small sample size of renal allograft biopsies with significant pathological findings. While 46 of 64 patients recruited for this study have undergone at least one biopsy, 36 of them were limited to zero-hour biopsies that did not yield any pathological findings. Furthermore, of the few indication biopsies performed, only 6 out of 39 combinations of lesion types and scores (\emph{ci}1, \emph{ct}1, and \emph{cv}1-\emph{cv}3 scores), as well as the IFTA1 grade, were encountered in 7 or more patients -- and only four patients experienced graft rejection during the study period. This limited the number of VisR measurements that could be associated with biopsy findings, particularly for higher \emph{i}, IFTA, and \emph{i-IFTA} scores as well as graft rejection. As a consequence, this study was unable to perform the longitudinal analyses that were originally envisioned, so only VisR measurements associated with biopsies were analyzed. This limitation raises two issues. First, the timing of VisR measurements relative to biopsies was not tightly controlled; patients were imaged whenever they presented for clinic visits regardless of when biopsies were scheduled. As such, there was considerable variability in the time elapsed between VisR measurements and biopsies (sometimes up to two months, especially during the COVID-19 pandemic), meaning that pathological changes may have occurred in the interim that were not captured by either modality. 

The inconsistency in the timing of VisR imaging sessions are especially noticeable beyond the first imaging session of each patient. Major practical limitations that affected patient recruitment and imaging frequency include the willingness of patients to attend follow-up visits in UNC's primary transplant nephrology clinic as opposed to satellite clinics or virtual visits (particularly during the COVID-19 pandemic), as well as their willingness to stay for an additional 15--20 minutes after their scheduled clinic visits for VisR imaging. These factors greatly affected the number of imaging sessions that could be performed for each patient. This hindered the ability to perform longitudinal analyses of VisR measurements relative to biopsy findings over time, especially for analyses of early detection where changes in VisR-derived DoA and RR metrics may precede pathological findings in biopsies. Future work should include more indication biopsies with significant pathological findings, as well as more consistent timing of VisR imaging sessions relative to biopsies, to investigate these questions.

This study also considered the detection of pathologies such as interstitial inflammation, IFTA, and graft rejection without fully separating out those populations. For example, this study observed the initial possibility of using VisR to detect IFTA \emph{in vivo} in a clinical setting through both DoA and RR techniques using PD, RE, and RV metrics -- but it did not investigate whether VisR could differentiate between IFTA and \emph{i-IFTA}, as biopsies could not be dichotomized in this manner. Considering that multiple pathologies may coexist in a single renal allograft, future work should investigate the ability of VisR-derived DoA and RR metrics to distinguish between different pathological conditions in renal allografts.

% Opted to not proceed with below paragraphs:
%
%Switching from methodological concerns to statistical ones, one must take care in interpreting the results from Tables~\ref{tbl:clinicalClassifier} and~\ref{tbl:clinicalRegressor}. While many of the clinical variables considered were categorical in nature, some metrics -- particularly IFTA grading and \emph{i-IFTA} scores -- were multiclass, ordered variables, requiring a regression model that does not inappropriately treat Banff scores as continuous variables. For example, it is not necessarily valid to grade a DoA that is ambiguously between IFTA1 and IFTA2 as "IFTA1.5", as this implies that the difference between IFTA1 and IFTA2 is the same as that between IFTA0 and IFTA1. The multinomial regression model used in this study addressed this concern by creating separate decision boundaries between each adjacent class (e.g. by ruling between IFTA0 and IFTA1 or above, then IFTA1 versus IFTA2 or above, etc.) such that the ordinal nature of variables are procedurally encoded but assumptions about the independence of irrelevant alternatives are avoided~\cite{long_Regression_2014}. This allowed the assessment of the correctness of model predictions (i.e. log-likelihood of predicting the correct class) as a distinct metric from the correlation strength. However, this also means that the AUC, specificity, and sensitivity metrics reported in Table~\ref{tbl:clinicalClassifier} represent the average performance of multiple binary classifiers rather than a single classifier per clinical variable. As such, it is difficult to generalize...
%
% AUC, specificity, and sensitivity thresholds of 0.5 corresponds to random guessing where all classes are equally likely. This resulted in a conservative estimate of model performance for renal allografts with pathologies, as they are less prevalent in the patient cohort for this study. Nevertheless, the models' performances were modest at best, with AUC values rarely exceeding 0.7 across all cross-validation folds. Likewise, the observed distributions of correlation coefficients and regressor error metrics exhibited relatively large standard deviations, indicating that some cross-validation folds yielded better performance than others.

VisR, as a imaging modality, also faces limitations in the context of this work. This study did not explicitly control for arterial occlusion, which has been shown to decrease renal stiffness along the dominant direction of nephrons~\cite{gennisson_Supersonic_2012,urban_Novel_2021,lim_Shear_2021}. Although the collagen deposition from IFTA may be confounded by perfusion-related changes in stiffness, the relative consistency of DoA and RR calculated from RE versus RV given an IFTA grade suggests that perfusion-related confounding effects may be limited in this study. A larger sample size, as well as direct measurements of renal blood flow using Doppler ultrasound or other imaging modalities, may help to clarify this relationship in future work.

Other practical and technique-related limitations must also be considered for VisR imaging. It is known that operator-dependent factors such as transducer placement and externally applied pre-compression can affect elastography parameters including shear wave-based techniques~\cite{barr_Effects_2012} as well as VisR~\cite{phillips_Surface_2025}. While this study attempted to mitigate these effects by using a single operator for all imaging sessions and providing standardized training on transducer placement and handling, some variability may still exist in the data due to these factors. Furthermore, patient-specific factors also limited the ability to obtain high-quality VisR measurements in certain cases, especially in the intended parts of the renal parenchyma, expected nephron orientation, or transducer orientation. While this was partly mitigated by using B-mode images with ARFI PD and log(VoA) overlays to guide ROI placements, which included the option for excluding images where ROIs could not be placed appropriately, some variability may still exist in the data due to these factors. Future work should investigate the effects of operator-dependent and patient-specific factors on VisR-derived DoA and RR metrics in renal allografts.

Lastly, VisR's use of relative metrics of elasticity and viscosity based on a simplified 1-D mass-spring-damper model of tissue response to ARF excitations should also be seen as an inherent limitation. This approach precluded the absolute measurement of elasticity and viscosity; DoA and RR were assumed to reflect relative differences in viscoelastic anisotropy and heterogeneity, respectively under the expectation that ARF push amplitudes were homogeneous across different acquisitions. If the spatial distribution of ARF push and track beams are distorted by aberration or focal mismatches, the accuracy of VisR-derived DoA and RR metrics may be compromised. While previous works~\cite{hossain_Viscoelastic_2018,hossain_Quantitative_2022} have demonstrated the consistency of VisR-derived DoA and RR metrics to finite element models (FEM), ground truth measurements, and other elastography techniques in tissue-mimicking phantoms and in animal models, further investigation specifically in humans is warranted to understand the stability of these metrics under more challenging conditions. A data-driven approaches to directly estimate shear elasticity or viscosity from ARF-induced displacements~\cite{richardson_Quantitative_2024} may help to address these concerns in future work.
